% preamble.tex -- pengaturan bersama seluruh gambar TikZ.
% Dipanggil dengan % preamble.tex -- pengaturan bersama seluruh gambar TikZ.
% Dipanggil dengan % preamble.tex -- pengaturan bersama seluruh gambar TikZ.
% Dipanggil dengan % preamble.tex -- pengaturan bersama seluruh gambar TikZ.
% Dipanggil dengan \input{preamble} tepat setelah \documentclass{standalone}.
%
% Font disamakan dengan template.cls (newtxtext/newtxmath) supaya teks di
% dalam gambar menyatu dengan teks naskah. Tanpa ini, angka pada sumbu grafik
% terlihat berbeda dari angka pada paragraf di sebelahnya.

\usepackage[T1]{fontenc}
\usepackage{newtxtext,newtxmath}
\usepackage{tikz}
\usepackage{pgfplots}
\pgfplotsset{compat=1.18}
\usetikzlibrary{calc,positioning,arrows.meta,decorations.pathreplacing,fit,backgrounds,shapes.geometric}

% Warna konsisten antar gambar.
\definecolor{aksen}{RGB}{0,90,160}      % biru: kurva atau objek utama kedua
\definecolor{aksendua}{RGB}{170,60,20}  % jingga gelap: objek ketiga
\definecolor{netral}{gray}{0.45}

% Gaya baku untuk diagram alur.
\tikzset{
  tahap/.style={
    rectangle, rounded corners=2pt, draw=black!55, fill=gray!8,
    align=center, font=\footnotesize, inner sep=5pt,
  },
  tahapaksen/.style={
    tahap, draw=aksen, fill=aksen!8,
  },
  panah/.style={
    -{Latex[length=5pt,width=4pt]}, line width=0.7pt, draw=black!60,
  },
  panahaksen/.style={
    panah, draw=aksen,
  },
  kelompok/.style={
    rectangle, rounded corners=3pt, draw=netral, dashed, line width=0.5pt,
    inner sep=7pt,
  },
}

% Gaya sumbu baku: garis tengah, grid tipis, label kecil.
\pgfplotsset{
  bakuaxis/.style={
    axis lines=left,
    axis line style={gray!70},
    tick label style={font=\footnotesize},
    label style={font=\small},
    legend style={font=\footnotesize, draw=gray!50, fill=white,
                  fill opacity=0.9, text opacity=1, row sep=1pt},
    grid=both,
    grid style={line width=0.2pt, draw=gray!18},
    every axis plot/.append style={line width=1.0pt},
  },
  bakubar/.style={
    ybar,
    bar width=\barw,
    fill=gray!35,
    draw=black!70,
    line width=0.5pt,
  },
}

\newcommand{\barw}{7pt}

% Label panel (a)/(b)/(c) di bawah sebuah axis bernama.
% Anchor "below south" milik pgfplots berada di bawah seluruh elemen axis,
% termasuk label sumbu, sehingga label panel tidak pernah bertabrakan dengannya.
\newcommand{\panellabel}[2]{%
  \node[font=\small, anchor=north, yshift=-1.5mm] at (#1.below south) {#2};}

% Keterangan bersama di bawah dua axis bernama.
\newcommand{\keteranganbersama}[3]{%
  \node[font=\scriptsize, align=center, anchor=north, yshift=-10mm]
    at ($(#1.below south)!0.5!(#2.below south)$) {#3};}
 tepat setelah \documentclass{standalone}.
%
% Font disamakan dengan template.cls (newtxtext/newtxmath) supaya teks di
% dalam gambar menyatu dengan teks naskah. Tanpa ini, angka pada sumbu grafik
% terlihat berbeda dari angka pada paragraf di sebelahnya.

\usepackage[T1]{fontenc}
\usepackage{newtxtext,newtxmath}
\usepackage{tikz}
\usepackage{pgfplots}
\pgfplotsset{compat=1.18}
\usetikzlibrary{calc,positioning,arrows.meta,decorations.pathreplacing,fit,backgrounds,shapes.geometric}

% Warna konsisten antar gambar.
\definecolor{aksen}{RGB}{0,90,160}      % biru: kurva atau objek utama kedua
\definecolor{aksendua}{RGB}{170,60,20}  % jingga gelap: objek ketiga
\definecolor{netral}{gray}{0.45}

% Gaya baku untuk diagram alur.
\tikzset{
  tahap/.style={
    rectangle, rounded corners=2pt, draw=black!55, fill=gray!8,
    align=center, font=\footnotesize, inner sep=5pt,
  },
  tahapaksen/.style={
    tahap, draw=aksen, fill=aksen!8,
  },
  panah/.style={
    -{Latex[length=5pt,width=4pt]}, line width=0.7pt, draw=black!60,
  },
  panahaksen/.style={
    panah, draw=aksen,
  },
  kelompok/.style={
    rectangle, rounded corners=3pt, draw=netral, dashed, line width=0.5pt,
    inner sep=7pt,
  },
}

% Gaya sumbu baku: garis tengah, grid tipis, label kecil.
\pgfplotsset{
  bakuaxis/.style={
    axis lines=left,
    axis line style={gray!70},
    tick label style={font=\footnotesize},
    label style={font=\small},
    legend style={font=\footnotesize, draw=gray!50, fill=white,
                  fill opacity=0.9, text opacity=1, row sep=1pt},
    grid=both,
    grid style={line width=0.2pt, draw=gray!18},
    every axis plot/.append style={line width=1.0pt},
  },
  bakubar/.style={
    ybar,
    bar width=\barw,
    fill=gray!35,
    draw=black!70,
    line width=0.5pt,
  },
}

\newcommand{\barw}{7pt}

% Label panel (a)/(b)/(c) di bawah sebuah axis bernama.
% Anchor "below south" milik pgfplots berada di bawah seluruh elemen axis,
% termasuk label sumbu, sehingga label panel tidak pernah bertabrakan dengannya.
\newcommand{\panellabel}[2]{%
  \node[font=\small, anchor=north, yshift=-1.5mm] at (#1.below south) {#2};}

% Keterangan bersama di bawah dua axis bernama.
\newcommand{\keteranganbersama}[3]{%
  \node[font=\scriptsize, align=center, anchor=north, yshift=-10mm]
    at ($(#1.below south)!0.5!(#2.below south)$) {#3};}
 tepat setelah \documentclass{standalone}.
%
% Font disamakan dengan template.cls (newtxtext/newtxmath) supaya teks di
% dalam gambar menyatu dengan teks naskah. Tanpa ini, angka pada sumbu grafik
% terlihat berbeda dari angka pada paragraf di sebelahnya.

\usepackage[T1]{fontenc}
\usepackage{newtxtext,newtxmath}
\usepackage{tikz}
\usepackage{pgfplots}
\pgfplotsset{compat=1.18}
\usetikzlibrary{calc,positioning,arrows.meta,decorations.pathreplacing,fit,backgrounds,shapes.geometric}

% Warna konsisten antar gambar.
\definecolor{aksen}{RGB}{0,90,160}      % biru: kurva atau objek utama kedua
\definecolor{aksendua}{RGB}{170,60,20}  % jingga gelap: objek ketiga
\definecolor{netral}{gray}{0.45}

% Gaya baku untuk diagram alur.
\tikzset{
  tahap/.style={
    rectangle, rounded corners=2pt, draw=black!55, fill=gray!8,
    align=center, font=\footnotesize, inner sep=5pt,
  },
  tahapaksen/.style={
    tahap, draw=aksen, fill=aksen!8,
  },
  panah/.style={
    -{Latex[length=5pt,width=4pt]}, line width=0.7pt, draw=black!60,
  },
  panahaksen/.style={
    panah, draw=aksen,
  },
  kelompok/.style={
    rectangle, rounded corners=3pt, draw=netral, dashed, line width=0.5pt,
    inner sep=7pt,
  },
}

% Gaya sumbu baku: garis tengah, grid tipis, label kecil.
\pgfplotsset{
  bakuaxis/.style={
    axis lines=left,
    axis line style={gray!70},
    tick label style={font=\footnotesize},
    label style={font=\small},
    legend style={font=\footnotesize, draw=gray!50, fill=white,
                  fill opacity=0.9, text opacity=1, row sep=1pt},
    grid=both,
    grid style={line width=0.2pt, draw=gray!18},
    every axis plot/.append style={line width=1.0pt},
  },
  bakubar/.style={
    ybar,
    bar width=\barw,
    fill=gray!35,
    draw=black!70,
    line width=0.5pt,
  },
}

\newcommand{\barw}{7pt}

% Label panel (a)/(b)/(c) di bawah sebuah axis bernama.
% Anchor "below south" milik pgfplots berada di bawah seluruh elemen axis,
% termasuk label sumbu, sehingga label panel tidak pernah bertabrakan dengannya.
\newcommand{\panellabel}[2]{%
  \node[font=\small, anchor=north, yshift=-1.5mm] at (#1.below south) {#2};}

% Keterangan bersama di bawah dua axis bernama.
\newcommand{\keteranganbersama}[3]{%
  \node[font=\scriptsize, align=center, anchor=north, yshift=-10mm]
    at ($(#1.below south)!0.5!(#2.below south)$) {#3};}
 tepat setelah \documentclass{standalone}.
%
% Font disamakan dengan template.cls (newtxtext/newtxmath) supaya teks di
% dalam gambar menyatu dengan teks naskah. Tanpa ini, angka pada sumbu grafik
% terlihat berbeda dari angka pada paragraf di sebelahnya.

\usepackage[T1]{fontenc}
\usepackage{newtxtext,newtxmath}
\usepackage{tikz}
\usepackage{pgfplots}
\pgfplotsset{compat=1.18}
\usetikzlibrary{calc,positioning,arrows.meta,decorations.pathreplacing,fit,backgrounds,shapes.geometric}

% Warna konsisten antar gambar.
\definecolor{aksen}{RGB}{0,90,160}      % biru: kurva atau objek utama kedua
\definecolor{aksendua}{RGB}{170,60,20}  % jingga gelap: objek ketiga
\definecolor{netral}{gray}{0.45}

% Gaya baku untuk diagram alur.
\tikzset{
  tahap/.style={
    rectangle, rounded corners=2pt, draw=black!55, fill=gray!8,
    align=center, font=\footnotesize, inner sep=5pt,
  },
  tahapaksen/.style={
    tahap, draw=aksen, fill=aksen!8,
  },
  panah/.style={
    -{Latex[length=5pt,width=4pt]}, line width=0.7pt, draw=black!60,
  },
  panahaksen/.style={
    panah, draw=aksen,
  },
  kelompok/.style={
    rectangle, rounded corners=3pt, draw=netral, dashed, line width=0.5pt,
    inner sep=7pt,
  },
}

% Gaya sumbu baku: garis tengah, grid tipis, label kecil.
\pgfplotsset{
  bakuaxis/.style={
    axis lines=left,
    axis line style={gray!70},
    tick label style={font=\footnotesize},
    label style={font=\small},
    legend style={font=\footnotesize, draw=gray!50, fill=white,
                  fill opacity=0.9, text opacity=1, row sep=1pt},
    grid=both,
    grid style={line width=0.2pt, draw=gray!18},
    every axis plot/.append style={line width=1.0pt},
  },
  bakubar/.style={
    ybar,
    bar width=\barw,
    fill=gray!35,
    draw=black!70,
    line width=0.5pt,
  },
}

\newcommand{\barw}{7pt}

% Label panel (a)/(b)/(c) di bawah sebuah axis bernama.
% Anchor "below south" milik pgfplots berada di bawah seluruh elemen axis,
% termasuk label sumbu, sehingga label panel tidak pernah bertabrakan dengannya.
\newcommand{\panellabel}[2]{%
  \node[font=\small, anchor=north, yshift=-1.5mm] at (#1.below south) {#2};}

% Keterangan bersama di bawah dua axis bernama.
\newcommand{\keteranganbersama}[3]{%
  \node[font=\scriptsize, align=center, anchor=north, yshift=-10mm]
    at ($(#1.below south)!0.5!(#2.below south)$) {#3};}
